\chapter{Cahn--Hilliard demo}
\label{chap:dev:cahn-hilliard}

\begin{quote}
Phase metadata. Spec dir:
\passthrough{\lstinline!specs/2026-05-03-phase-12-cahn-hilliard/!}.
First scientifically-meaningful end-to-end demo: Cahn--Hilliard
coarsening on a 3D periodic grid with explicit RK4 and a hand-written
legacy-VTK snapshot writer. Version target:
\passthrough{\lstinline!0.12.0!}.
\end{quote}

\subsection{1. Theory}\label{phase12-theory}

\subsubsection{The Cahn--Hilliard equation}\label{phase12-ch}

Cahn \& Hilliard (1958) introduced the conservative gradient flow
\[
\partial_t\psi(\mathbf x, t) \;=\; M\,\nabla^2\!\Big(\frac{\delta F}{\delta \psi}\Big),
\qquad
F[\psi] \;=\; \int_\Omega f(\psi, \nabla\psi)\,\dd V,
\]
as a model of phase separation in a binary alloy. The order parameter
\(\psi\) is the local composition (or its deviation from a reference
mean); \(M>0\) is the mobility. The dynamics is conservative ---
\(\partial_t\!\int\psi = \oint(\text{flux})\cdot\dd\mathbf S = 0\) on a
periodic or no-flux domain --- and is a gradient descent of \(F\) in
the \(H^{-1}\) Sobolev metric induced by the
\passthrough{\lstinline!\nabla^2!} prefactor.

The standard Landau / Ginzburg form of the free-energy density,
\[
f(\psi, \nabla\psi) \;=\; -\tfrac12\psi^2 + \tfrac14\psi^4 + \tfrac{\kappa}{2}|\nabla\psi|^2,
\]
combines a non-convex bulk term with double-well minima at
\(\psi=\pm 1\) and the standard interfacial-gradient penalty. The
parameter \(\kappa>0\) sets the interfacial width
\(\xi\sim\sqrt{\kappa}\) and the surface tension
\(\sigma\sim\sqrt{\kappa}\) of the equilibrium domain wall (Cahn 1958).

\subsubsection{Lyapunov property}\label{phase12-lyapunov}

Multiplying the CH equation by
\(\mu := \delta F/\delta\psi\) and integrating by parts:
\[
\frac{\dd F}{\dd t}
   \;=\; \int \mu\,\partial_t\psi\,\dd V
   \;=\; \int \mu\,M\nabla^2\mu\,\dd V
   \;=\; -M\!\int |\nabla\mu|^2\,\dd V \;\le\; 0.
\]
\(F\) is therefore a Lyapunov functional --- the dynamics monotonically
decreases \(F\) until \(\nabla\mu\equiv 0\), i.e., until \(\mu\) is
spatially uniform (the equilibrium condition for two-phase coexistence
at fixed mean composition). The Phase 12 acceptance test
\passthrough{\lstinline!CahnHilliardCoarseningRun.LyapunovDecay!}
locks this property in.

\subsubsection{Spinodal decomposition and the
coarsening regime}\label{phase12-coarsening}

Around the unstable mean \(\bar\psi=0\), the linearized dynamics is
\[
\partial_t\delta\psi \;=\; M\,\nabla^2(-\delta\psi - \kappa\nabla^2\delta\psi)
\quad\Longrightarrow\quad
\sigma(\mathbf k) \;=\; M\,|\mathbf k|^2\,(1 - \kappa|\mathbf k|^2).
\]
\(\sigma>0\) for \(|\mathbf k|<\sqrt{\kappa^{-1}}\) (the spinodal
band), maximized at
\(k_*=1/\sqrt{2\kappa}\) with growth rate
\(\sigma_{\max}=M/(4\kappa)\). Random ICs of small amplitude grow
exponentially in the \(k_*\) shell until the cubic \(\psi^3\) term
saturates at \(|\psi|\approx 1\).

After saturation, the system enters the classical-coarsening regime
described by Bray (1994): the typical domain size \(R(t)\) grows as
\(t^{1/3}\) (Lifshitz--Slyozov scaling for a conserved scalar order
parameter). The interfacial-area density and the structure-factor
peak position \(k_*(t)\) both decrease as \(t^{-1/3}\). At the demo's
finite size, the coarsening regime is captured qualitatively: the
gradient energy \(E_{\text{grad}}\) peaks near
\(k_*\approx k_{*,\text{linear}}\), then decreases monotonically.

\subsection{2. Math derivation}\label{phase12-math}

\subsubsection{Closed-form variational
derivative}\label{phase12-deriv}

For the Landau-double-well + gradient form,
\(\delta F/\delta\psi\) is computed by the standard
Euler--Lagrange recipe (taking
\(\partial f/\partial\psi - \partial_a (\partial f/\partial\psi_{,a})\)):
\begin{align*}
\frac{\delta F}{\delta\psi}
&\;=\; \frac{\partial f}{\partial\psi} - \nabla\!\cdot\!\frac{\partial f}{\partial(\nabla\psi)} \\
&\;=\; (-\psi + \psi^3) - \nabla\!\cdot\!(\kappa\nabla\psi) \\
&\;=\; -\psi + \psi^3 - \kappa\nabla^2\psi.
\end{align*}
The full RHS is \(M\nabla^2\) of this: two compositions of
\passthrough{\lstinline!diff::laplacian!} per RHS evaluation, with a
pointwise \(-\psi+\psi^3\) sandwiched between them.

Phase 12 supplies this in closed form (no AD): the helper
\passthrough{\lstinline!gridcalc::examples::cahn\_hilliard::computeRhs!}
is a verbatim transcription of
\(M\,\nabla^2(-\psi + \psi^3 - \kappa\nabla^2\psi)\) and is the only
non-trivial mathematical content of
\href{../../../examples/common/cahn_hilliard.hpp}{\passthrough{\lstinline!examples/common/cahn\_hilliard.hpp!}}.

\subsubsection{Linear stability and the explicit-RK4 step
budget}\label{phase12-rk4}

The discrete operator \(\diff::laplacian\) has eigenvalues on the
periodic grid \(\Lambda \in [-12/\Delta x^2,\,0]\) (each axis
contributes \([-4/\Delta x^2,\,0]\) and they sum). Composed twice for
the linearized RHS \(-M\nabla^2\psi - M\kappa\nabla^4\psi\), the
eigenvalues become
\[
\omega(\Lambda) \;=\; -M\,\Lambda - M\,\kappa\,\Lambda^2 \;=\; -M\Lambda\,(1 + \kappa\Lambda).
\]
For \(\Lambda<0\) and \(|\kappa\Lambda|>1\), the second term
dominates: \(\omega \approx -M\,\kappa\,\Lambda^2\).
The most negative discrete eigenvalue is therefore at
\(\Lambda_{\min} = -12/\Delta x^2\):
\[
|\omega|_{\max} \;\approx\; M\,\kappa\,\frac{144}{\Delta x^4}.
\]

Classical RK4 has a stability region intersecting the negative real
axis at \(\Re(z)\in [-2.7853, 0]\) (LeVeque 2007, Chapter 7;
Hairer--Nørsett--Wanner 1993). Stable explicit time stepping requires
\(\Delta t\,|\omega|_{\max}\le 2.7853\), giving
\[
\Delta t_{\max} \;=\; \frac{2.7853\,\Delta x^4}{144\,M\,\kappa}.
\]

For the demo defaults \(M=\kappa=1\), \(\Delta x=1\):
\(\Delta t_{\max}\approx 0.0193\). The default
\passthrough{\lstinline!--dt 0.01!} leaves a 1.93\(\times\) safety
factor; the test parameters \(\kappa=0.5\), \(\Delta t=0.02\) leave
1.93\(\times\) as well (\(\Delta t_{\max} \approx 0.0387\) at
\(\kappa=0.5\)).

The \(\Delta x^4\) scaling is the central pain point: doubling the
resolution divides the allowed step by 16. This is the structural
reason implicit / IMEX schemes are essential for production-scale CH
simulations and why Phase 19 is on the roadmap.

\subsubsection{Why no \texttt{laplacian}\(\circ\)\texttt{laplacian}}\label{phase12-no-direct}

The CH RHS could be written as
\(M\,\nabla^2(-\psi + \psi^3) - M\kappa\,\nabla^2(\nabla^2\psi)\), with
the second term coming from
\passthrough{\lstinline!laplacian(laplacian(psi))!}. We instead use
\(M\,\nabla^2(\mu)\) where
\(\mu = -\psi + \psi^3 - \kappa\nabla^2\psi\) is computed once
explicitly. The pointwise \(\psi^3\) factor in \(\mu\) couples the
high-frequency content of \(\psi\) into the outer Laplacian's
stencil radius, which is a feature, not a bug --- it captures the
correct nonlinear contribution at the discrete level. Splitting the
two operators and applying the second axis-decomposed Laplacian to
\((\psi - \psi^3)\) separately would produce a different
discretization with different leading-order error. The closed-form
recipe is unambiguous; we follow it.

\subsection{3. Algorithm}\label{phase12-algorithm}

\subsubsection{File layout}\label{phase12-files}

\begin{itemize}
\tightlist
\item
  \href{../../../examples/cahn_hilliard.cpp}{\passthrough{\lstinline!examples/cahn\_hilliard.cpp!}}
  --- \passthrough{\lstinline!main()!}, CLI parsing, RK4 time-stepping
  loop, per-snapshot diagnostics + VTK write.
\item
  \href{../../../examples/common/cahn_hilliard.hpp}{\passthrough{\lstinline!examples/common/cahn\_hilliard.hpp!}}
  --- \passthrough{\lstinline!computeRhs!},
  \passthrough{\lstinline!computeFreeEnergy!},
  \passthrough{\lstinline!computeGradientEnergy!},
  \passthrough{\lstinline!computeMean!},
  \passthrough{\lstinline!makeRandomIc!}.
\item
  \href{../../../examples/common/vtk_writer.hpp}{\passthrough{\lstinline!examples/common/vtk\_writer.hpp!}}
  --- \passthrough{\lstinline!writeVtkStructuredPoints!} (legacy ASCII
  VTK 3.0 \passthrough{\lstinline!STRUCTURED\_POINTS!} writer).
\item
  \href{../../../test/cahn_hilliard_test.cpp}{\passthrough{\lstinline!test/cahn\_hilliard\_test.cpp!}}
  --- six tests including the three TEST\_F coarsening-property
  assertions on a single shared 30\,s simulation.
\item
  \href{../../../test/vtk_writer_test.cpp}{\passthrough{\lstinline!test/vtk\_writer\_test.cpp!}}
  --- three round-trip checks of the VTK writer (header shape,
  payload exact-match in i-fastest order, scalars-line field-name
  round-trip).
\item
  \href{../../../scripts/render_ch_snapshots.py}{\passthrough{\lstinline!scripts/render\_ch\_snapshots.py!}}
  --- matplotlib helper that renders 2D z-midplane PNGs from the
  written \passthrough{\lstinline!.vtk!} snapshots.
\end{itemize}

\subsubsection{Build wiring}\label{phase12-build}

A new top-level option
\passthrough{\lstinline!GRIDCALC\_BUILD\_EXAMPLES!} (default
\passthrough{\lstinline!OFF!}, mirroring Phase 10's
\passthrough{\lstinline!GRIDCALC\_BUILD\_DOCS!}) gates an
\passthrough{\lstinline!add\_subdirectory(examples)!} from the root
\passthrough{\lstinline!CMakeLists.txt!}. The child file
\passthrough{\lstinline!examples/CMakeLists.txt!} declares the
\passthrough{\lstinline!cahn\_hilliard!} target, links it against
\passthrough{\lstinline!gridcalc::gridcalc!} and
\passthrough{\lstinline!gridcalc\_warnings!}, and adds
\passthrough{\lstinline!examples/!} to its include path so the helper
headers under \passthrough{\lstinline!examples/common/!} are reachable
as \passthrough{\lstinline!\#include <common/...>!}.

CI (both Ubuntu GCC and Ubuntu Clang) configures with
\passthrough{\lstinline!-DGRIDCALC\_BUILD\_EXAMPLES=ON!} so the
example is type-checked under the project's strict warning set on
every PR. The example binary is \emph{not} run in CI; the runtime
gate is the acceptance test in \passthrough{\lstinline!test/!}, which
includes the helpers as headers and is therefore active regardless
of \passthrough{\lstinline!GRIDCALC\_BUILD\_EXAMPLES!}.

\subsubsection{Time-stepping loop and per-step
allocations}\label{phase12-loop}

Time stepping is \passthrough{\lstinline!solve::integrate(psi, rhs,
dt, n, RK4{})!} (Phase 6) applied in chunks of
\passthrough{\lstinline!--snapshot-every!} steps. Per RK4 step:
\begin{itemize}
\tightlist
\item \(4\times\) RHS evaluations
   (\passthrough{\lstinline!computeRhs!}),
\item each \passthrough{\lstinline!computeRhs!} allocates 3 fresh
   \passthrough{\lstinline!Field<double>!}: one for
   \passthrough{\lstinline!lap\_psi!}, one for the chemical potential
   \(\mu\), one for the outer \passthrough{\lstinline!lap\_mu!} which
   is the returned RHS,
\item plus the four
   \passthrough{\lstinline!detail::axpyFresh!}
   intermediate fields (\(y_2\), \(y_3\), \(y_4\), and the four
   \(k_i\)) that Phase 6's RK4 implementation already materializes,
\item per step: \(4 \cdot 3 + 8 = 20\)
   \passthrough{\lstinline!Field<double>!} allocations.
\end{itemize}

At \(N=64\) (\(64^3 = 262{,}144\) doubles \(\approx 2.1\,\)MB per
\passthrough{\lstinline!Field<double>!}), the per-step transient
working set is \(\sim 42\,\)MB --- well within working-set targets but
representative of the per-step churn that Phase 20's performance pass
will replace with a persistent scratch pool.

\subsubsection{Snapshot writing}\label{phase12-snapshot}

\passthrough{\lstinline!writeVtkStructuredPoints!} writes the legacy
ASCII VTK 3.0 format:

\begin{lstlisting}
# vtk DataFile Version 3.0
gridcalc field: psi
ASCII
DATASET STRUCTURED_POINTS
DIMENSIONS Nx Ny Nz
ORIGIN 0 0 0
SPACING hx hy hz
POINT_DATA Nx*Ny*Nz
SCALARS psi double 1
LOOKUP_TABLE default
<value at flat index 0>
<value at flat index 1>
...
\end{lstlisting}

The data block is a flat read of \passthrough{\lstinline!field.data()!}
in i-fastest order (\(\text{idx} = i + N_x j + N_x N_y k\)) --- exactly
the order VTK \passthrough{\lstinline!STRUCTURED\_POINTS!} expects, so
no transposition is needed. Values are emitted with 17-digit
decimal precision (\passthrough{\lstinline!std::ofstream::precision(17)!})
to round-trip a \passthrough{\lstinline!double!} losslessly.

ASCII (rather than binary or XML \passthrough{\lstinline!.vti!}) is
the deliberate choice: \(\sim 30\) lines of hand-written I/O, no new
build dependency, human-readable for debugging. The file size is
\(\sim 25\,\)bytes per grid point in ASCII vs.\ \(8\) for binary; at
\(N=64\) that is \(6.5\,\)MB per snapshot --- inconvenient at very
high resolution but irrelevant for the demo budget. Production users
needing binary or XML VTK should integrate
\passthrough{\lstinline!vtkIO!} directly.

\subsubsection{Acceptance-test instrumentation}\label{phase12-acc}

\passthrough{\lstinline!test/cahn\_hilliard\_test.cpp!} runs a small
CH simulation once per process inside a static-cached lambda
\passthrough{\lstinline!Snaps()!}. Three \passthrough{\lstinline!TEST\_F!}
bodies share the result, so the \(\sim 30\,\)s simulation does not
repeat per test:

\begin{lstlisting}[language={C++}]
static const std::vector<Snapshot>& Snaps() {
    static const std::vector<Snapshot> snaps = []() {
        // ...build IC, run 4 RK4 segments, collect 5 snapshots...
    }();
    return snaps;
}
\end{lstlisting}

The test parameters
(\(N=24\), \(\Delta t=0.02\), \(\kappa=0.5\), \(M=1\), amplitude
\(0.1\), seed 42, 4 segments of \{500, 500, 500, 1000\} steps for
snapshots at \(t \in \{0,10,20,30,50\}\)) are calibrated so the
\(t=30\) snapshot lies just past the gradient-energy peak and the
\(t=50\) snapshot is one decade into the coarsening regime. Locally
measured gradient-energy drops across seeds \(\{1, 2, 7, 42, 1234\}\)
range from 4.83\,\% to 7.85\,\%; the test asserts a 3\,\% drop ---
robust against \passthrough{\lstinline!std::uniform\_real\_distribution!}
differences across libstdc++ and libc++.

\subsubsection{Snapshot rendering and figure
generation}\label{phase12-figures}

\passthrough{\lstinline!scripts/render\_ch\_snapshots.py!} parses the
ASCII VTK header positionally (the writer's output is fully
deterministic), reshapes the 1D payload back to \((N_z, N_y, N_x)\)
(the inverse of the i-fastest layout), and renders the z-midplane
slice with the \passthrough{\lstinline!matplotlib!} ``Agg'' backend
(no display dependency). The three figures committed under
\passthrough{\lstinline!docs/user-guide/figures/cahn-hilliard/!}
were generated from a \(64^3\) release-build run at the demo's
default parameters at steps \(\{2000, 4000, 8000\}\) (corresponding
to \(t \in \{20, 40, 80\}\)).

\subsection{4. Design decisions}\label{phase12-decisions}

These distill the choices recorded in
\passthrough{\lstinline!specs/2026-05-03-phase-12-cahn-hilliard/requirements.md!}.

\begin{itemize}
\item
  \textbf{Standard Landau-double-well + gradient term, hand-coded
  \(\delta F/\delta\psi\).} The textbook CH form, parameterless
  except for the explicit \(\kappa\) and \(M\) (both default to 1).
  Closed-form variational derivative; no automatic differentiation
  anywhere on the roadmap. Logarithmic / Flory--Huggins forms were
  considered and rejected for this phase --- they add
  \passthrough{\lstinline!(0,1)!}-clamping concerns that obscure the
  demo's qualitative-coarsening message.
\item
  \textbf{Explicit RK4 via Phase 6's
  \passthrough{\lstinline!solve::integrate!}.} Zero new solver code.
  Semi-implicit / IMEX paths (which would need
  Fourier-diagonalization, violating the
  \passthrough{\lstinline!mission.md!} ``FFT verification only''
  decision) and matrix-free implicit schemes (Phase 19's natural
  scope) are explicitly out of scope. The brutal
  \(\Delta t \sim h^4\) scaling is documented; resolution is small
  enough in the demo that the budget stays sub-minute.
\item
  \textbf{Legacy ASCII VTK \passthrough{\lstinline!STRUCTURED\_POINTS!}.}
  Hand-written, no new dependency, opens directly in ParaView and
  VisIt. Binary VTK was considered (\(5\)--\(10\times\) smaller files,
  endian-handling complexity); XML \passthrough{\lstinline!.vti!} was
  considered (modern format, more code to write correctly without a
  library). For a demo budget the ASCII form is correct.
\item
  \textbf{Per-step interval cadence (\passthrough{\lstinline!--snapshot-every!})
  rather than wall-clock.} The simulation is deterministic; clock
  reading would add machinery for no scientific benefit. The CLI
  exposes the cadence so the user can tune snapshot density without
  rebuilding.
\item
  \textbf{Fixed-seed mean-subtracted random IC.} The random-uniform
  amplitude is mean-subtracted to enforce \(\Sigma\psi=0\) at machine
  precision, ensuring CH's exact mass conservation starts from a
  zero-mean state. \passthrough{\lstinline!std::mt19937\_64!} is
  deterministic across implementations; only
  \passthrough{\lstinline!std::uniform\_real\_distribution!} differs.
  The acceptance test's
  \passthrough{\lstinline!SeedReproducibility!} check verifies
  in-implementation reproducibility.
\item
  \textbf{Visual + monotone-coarsening acceptance bar.} Two evidence
  streams: three committed PNG snapshots in the User Guide chapter
  (qualitative match to published CH images) plus the programmatic
  \passthrough{\lstinline!CahnHilliardCoarseningRun!} tests in CI
  (Lyapunov decay, monotone coarsening of the gradient energy, mass
  conservation). A quantitative \(R(t)\sim t^{1/3}\) power-law fit was
  considered and rejected for this phase: at the affordable grid size
  the slope estimate is noisy enough that any threshold band wide
  enough to be robust would also be uninformative. Could ship later
  as a separate non-CI benchmark.
\item
  \textbf{Helper headers under
  \passthrough{\lstinline!examples/common/!}, not in
  \passthrough{\lstinline!include/gridcalc/!}.} The CH RHS, energy
  diagnostics, IC generator, and VTK writer are not part of the
  library's public API. The Phase 10 \passthrough{\lstinline!\\since!}
  lint regex covers \passthrough{\lstinline!include/gridcalc/!}; placing
  the example helpers there would either force them onto the public
  API surface or require a lint-regex carve-out for a non-library
  subdirectory. Keeping them in
  \passthrough{\lstinline!examples/common/!} avoids both. Tests
  \passthrough{\lstinline!\#include!} the helpers via the
  \passthrough{\lstinline!target\_include\_directories!} entry added
  in \passthrough{\lstinline!test/CMakeLists.txt!}.
\end{itemize}

\subsection{5. References}\label{phase12-references}

External:

\begin{itemize}
\tightlist
\item
  \textbf{Cahn, J. W., \& Hilliard, J. E. (1958). \emph{Free Energy
  of a Nonuniform System. I. Interfacial Free Energy}. Journal of
  Chemical Physics, 28(2), 258--267.}
  \url{https://doi.org/10.1063/1.1744102}. Original derivation of the
  gradient-energy functional and its variational derivative; the
  source of the Landau-double-well + gradient form used here.
\item
  \textbf{Cahn, J. W. (1959). \emph{Free Energy of a Nonuniform
  System. II. Thermodynamic Basis}. Journal of Chemical Physics,
  30(5), 1121--1124.} \url{https://doi.org/10.1063/1.1730145}.
  Companion paper deriving the dynamical equation
  \(\partial_t\psi = M\nabla^2\mu\) on thermodynamic grounds.
\item
  \textbf{Bray, A. J. (1994). \emph{Theory of phase-ordering
  kinetics}. Advances in Physics, 43(3), 357--459.}
  \url{https://doi.org/10.1080/00018739400101505}. The canonical
  reference on classical-coarsening scaling (\(R(t)\sim t^{1/3}\) for
  conserved scalar order parameters; \(t^{1/2}\) for non-conserved);
  Section 4 covers Cahn--Hilliard specifically.
\item
  \textbf{Provatas, N., \& Elder, K. R. (2010). \emph{Phase-Field
  Methods in Materials Science and Engineering}. Wiley-VCH.} ISBN
  978-3-527-40747-7. Chapters 2--3 cover the Cahn--Hilliard
  derivation, equilibrium interface profile, and surface tension
  expressions used informally in
  \Cref{phase12-ch}.
\item
  \textbf{LeVeque, R. J. (2007). \emph{Finite Difference Methods for
  Ordinary and Partial Differential Equations}. SIAM.}
  \url{https://doi.org/10.1137/1.9780898717839}. Chapter 7 (``Absolute
  stability for ODEs'') gives the RK4 stability region and the
  \(2.7853\) intersection with the negative real axis used in
  \Cref{phase12-rk4}.
\item
  \textbf{Hairer, E., Nørsett, S. P., \& Wanner, G. (1993).
  \emph{Solving Ordinary Differential Equations I: Nonstiff
  Problems} (2nd ed.). Springer-Verlag.} ISBN 978-3-540-56670-0.
  Section II.4 derives the classical RK4 stability polynomial and
  region; cross-checked against the LeVeque value.
\item
  \textbf{VTK File Formats reference.} Permanent URL:
  \url{https://kitware.github.io/vtk-examples/site/VTKFileFormats/}.
  Defines the legacy \passthrough{\lstinline!STRUCTURED\_POINTS!}
  ASCII format used by \passthrough{\lstinline!writeVtkStructuredPoints!}.
\end{itemize}

Internal cross-references (do not satisfy the ``external'' minimum):

\begin{itemize}
\tightlist
\item
  \Cref{chap:dev:periodic-3d-scalar-laplacian} --- Phase 1's
  \passthrough{\lstinline!diff::laplacian!}, applied twice per CH RHS
  evaluation.
\item
  \Cref{chap:dev:rk4--generic-time-integrator} --- Phase 6's
  tag-dispatched \passthrough{\lstinline!solve::integrate(..., RK4{})!}
  used as the time-stepping driver.
\item
  \Cref{chap:dev:functional-evaluation} --- Phase 4's
  \passthrough{\lstinline!func::evaluate!}, used to compute \(F\),
  \(E_{\text{grad}}\), and \(\langle\psi\rangle\) over the demo state.
\end{itemize}
